\begin{proposition}[The bivariate moving-gap kernel]
\label{prop:gap-kernel-bivariate}
Let
\[
 B(x)=\sum_{n\geq0}b_nx^n,
 \qquad
 C(x)=\sum_{n\geq0}c_nx^n,
\]
where $(b_0,b_1)=(1,5)$ and $(c_0,c_1)=(0,1)$, and retain the
normalized restart coefficients $V_m(a)$ from
Proposition~\ref{prop:gap-kernel-green}.  Put
\[
 \mathcal G(s,z)=\sum_{a\geq1}\sum_{m\geq0}V_m(a)s^az^m
\]
and, with \(\theta_x=x\,d/dx\), define
\begin{equation}
 \mathcal K(X,Z)
 =X\left(
   (\theta_X+1)^3B(X)C(Z)
   -(\theta_X+1)^3C(X)B(Z)
  \right).
 \label{eq:bivariate-K}
\end{equation}
Then the following formal identities hold.

\begin{enumerate}
\item If \(\mathcal P_z^{\geq0}\) denotes extraction of the nonnegative
  powers of~$z$, then
  \begin{equation}
   \boxed{
   \mathcal G(s,z)
    =\mathcal P_z^{\geq0}\mathcal K(s/z,z)
    =\operatorname{CT}_u
      \frac{\mathcal K(su/z,z/u)}{1-u}.}
   \label{eq:bivariate-positive-part}
  \end{equation}
  The constant term is well defined coefficientwise in~$s$: for a fixed
  power~$s^a$, only finitely many negative powers of~$z$ occur before
  applying~$\mathcal P_z^{\geq0}$.

\item Let
  \[
   \Theta=\theta_s+\theta_z,
   \qquad
   \mathscr D_{s,z}
    =\Theta^3-zP(\Theta)+z^2(\Theta+1)^3.
  \]
  Then
  \begin{equation}
   \boxed{
   \mathscr D_{s,z}\mathcal G(s,z)
    =\sum_{a\geq1}a^3s^a
    =\frac{s(1+4s+s^2)}{(1-s)^4}.}
   \label{eq:bivariate-PDE}
  \end{equation}
  The two boundary slices are
  \begin{equation}
   [s]\mathcal G(s,z)=\frac{C(z)}z,
   \qquad
   [z^0]\mathcal G(s,z)=\frac{s}{1-s}.
   \label{eq:bivariate-boundary-slices}
  \end{equation}
\end{enumerate}
\end{proposition}

\begin{proof}
The companion generating functions satisfy
\begin{equation}
 \mathscr L_0B=0,
 \qquad
 \mathscr L_0C=z,
 \qquad
 \mathscr L_0=\theta^3-zP(\theta)+z^2(\theta+1)^3.
 \label{eq:companion-generating-equations}
\end{equation}
Indeed, the coefficients of degree at least two are exactly the Ap\'ery
recurrence, while the coefficient of~$z$ in the second equation is
$c_1-P(0)c_0=1$.

The coefficient of $X^{r+1}Z^n$ in~\eqref{eq:bivariate-K} is
\[
 (r+1)^3\bigl(b_rc_n-c_rb_n\bigr).
\]
After substituting $X=s/z$ and $Z=z$, this term becomes
\[
 (r+1)^3\bigl(b_rc_n-c_rb_n\bigr)
 s^{r+1}z^{n-r-1}.
\]
Its $z$-power is nonnegative exactly when $n\geq r+1$.  Setting
$a=r+1$ and $m=n-r-1$, the Casoratian formula
\eqref{eq:green-casoratian} identifies the coefficient with~$V_m(a)$.
This proves the first equality in~\eqref{eq:bivariate-positive-part}.
For a Laurent series $A(t)=\sum_{k\in\mathbb Z}A_kt^k$ one has
\[
 \operatorname{CT}_u\frac{A(z/u)}{1-u}
 =\sum_{k\geq0}A_kz^k.
\]
Applying this identity to $A(t)=\mathcal K(s/t,t)$ proves the constant-term
formula.

For each fixed~$a$, Proposition~\ref{prop:gap-kernel-green} gives
\[
 \bigl((\theta_z+a)^3-zP(\theta_z+a)
       +z^2(\theta_z+a+1)^3\bigr)F_a(z)=a^3.
\]
Multiplication by~$s^a$ followed by summation over $a\geq1$ replaces
$a$ by~$\theta_s$ and proves~\eqref{eq:bivariate-PDE}.  Finally,
$V_m(1)=c_{m+1}$ and $V_0(a)=1$, which give
\eqref{eq:bivariate-boundary-slices}.
\end{proof}

\begin{proposition}[Infinite-denominator obstruction]
\label{prop:gap-kernel-denominator-obstruction}
For every prime $p\geq5$,
\begin{equation}
 v_p(c_p)=-3,
 \qquad
 p^3c_p\equiv1\pmod p.
 \label{eq:companion-prime-denominator}
\end{equation}
Consequently neither $C(z)$ nor $\mathcal G(s,z)$ is globally bounded:
no fixed rescaling of its variables and of the whole series makes all its
coefficients integral.  In particular, $\mathcal G$ cannot be obtained as
an ordinary diagonal or constant term of a fixed rational function over
$\mathbb Q$ that is regular at the expansion point.
\end{proposition}

\begin{proof}
The values $c_0,\ldots,c_{p-1}$ are $p$-integral because every recurrence
denominator used to construct them is a unit modulo~$p$.  Reflection of a
$p$-integral solution gives
\[
 c_{p-1}\equiv c_0=0,
 \qquad
 c_{p-2}\equiv c_1=1
 \pmod p.
\]
The recurrence at index $p-1$ is
\[
 p^3c_p=P(p-1)c_{p-1}-(p-1)^3c_{p-2}.
\]
Reduction modulo~$p$ gives $p^3c_p\equiv1$, proving
\eqref{eq:companion-prime-denominator}.

By~\eqref{eq:bivariate-boundary-slices},
\[
 [s z^{p-1}]\mathcal G(s,z)=c_p.
\]
Thus infinitely many distinct primes occur in coefficient denominators of
both $C$ and~$\mathcal G$, and a fixed integral rescaling cannot remove
them.  On the other hand, after multiplying a rational function over
$\mathbb Q$ by one fixed integer, the coefficients of its expansion at a
regular point lie in $\mathbb Z[1/N]$ for one fixed~$N$.  Any diagonal or
constant-term extraction preserves this finite set of denominator primes.
The asserted rational-diagonal obstruction follows.
\end{proof}

\begin{remark}[Scope of the obstruction]
\label{rem:bivariate-kernel-scope}
Formula~\eqref{eq:bivariate-positive-part} compresses the starting index
and bridge length into one fixed formal construction, but it necessarily
uses the inhomogeneous companion series~$C$.  Proposition
\ref{prop:gap-kernel-denominator-obstruction} rules out the most direct
globally bounded rational-diagonal realization.  It does not rule out a
mixed-motive or regulator-valued realization, nor a characteristic-$p$
construction whose arithmetic model varies with~$p$.  In particular, the
constant-term formula alone supplies no uniform conductor bound and no
nonconcentration estimate for the coefficients~$V_m(a)$ modulo~$p$.
\end{remark}
